\scanpage{242}
- \mspage{182} - \margin{IX}
13.5. \uline{Cas d'une situation relativement algébrique}. Nous supposons maintenant que $A$ provient d'un schéma en groupes relatif sur $S$, lisse, séparé et à fibres des extensions de schémas abéliens par des tores. Nous désignerons ce schéma relatif par la même lettre $A$. (Nous nous autorisons cet abus de notation grâce au fait que le foncteur $G\longmapsto G^{\mathrm{an}}$, qui va de la catégorie des schémas en groupes relatifs du type envisagé vers la catégorie des groupes analytiques sur $S$, est pleinement fidèle, comme il serait facile à établir grâce aux résultats du type GAGA relatifs [10].) Pour tout ce qui concerne la notion de schéma relatif, et les relations entre schémas relatifs sur $S$ et espaces analytiques sur $S$, nous renvoyons à [10] (cf. aussi [2, Exp. 15]). Si $\underline O$ désigne l'anneau local de $S$ en $s$, le schéma relatif $A$ est connu au voisinage de $s$ par la connaissance du schéma ordinaire $A_{\underline O}$ sur $\operatorname{Spec}(\underline O)$, déduit de $A$ par localisation en $s$. Il est clair alors que le complété formel de $A_{\underline O}$ le long de sa fibre spéciale (isomorphe à $A_s=A_0$) est isomorphe au complété formel du groupe analytique $A$ le long de $A_s$ (13.2), en identifiant les $A_n=A\times_S S_n$ à des schémas algébriques sur $S$, par l'abus de notations introduit ci-dessus. Dans le cas favorable où $\widehat A$, regardé comme schéma formel en groupes sur $\widehat{\underline O}$, est algébrisable (7.2.1), ce qui est le cas en particulier si $S$ est normal en $s$ et $A$ quasi-projectif sur $S$ (7.1.5) on obtient donc un schéma en groupes $A_{\widehat{\underline O}}^{\unsure{\natural|\sharp}}$ sur $\operatorname{Spec}(\widehat{\underline O})$, extension d'un schéma abélien $B_{\widehat{\underline O}}$ par un tore. Pour simplifier, nous supposons que $B_{\widehat{\underline O}}$ est projectif sur $\widehat{\underline O}$ (ce qui est bien le cas, on le sait, si $\underline O$ donc $\widehat{\underline O}$ est normal). On va alors prouver :
\margin{\unsure{1}}
\margin{500}
\margin{241}
\margin{500}
